% Vorlage Diagramm (PGFPlots) im Leitfaden-Stil
%
% Uebersetzen mit LuaLaTeX. Im selben Ordner muessen liegen:
%   vfdbstil.sty, vfdbstil-werte.tex   (und ggf. die Datendatei)
% Ergebnis: PDF in fester Groesse -> in Overleaf nach media/ hochladen und
% ohne Skalierung einbinden: \includegraphics{media/XX-name.pdf}
%
% WICHTIG: Die Abbildung hat eine feste Groesse. Was ausserhalb liegt (z. B. eine
% Legendenzeile zu viel), wird ohne Warnung abgeschnitten - das PDF ansehen.
% SPDX-License-Identifier: MIT
\documentclass[border=0pt]{standalone}
\usepackage[ngerman]{babel}
\usepackage[schriften]{vfdbstil}

\begin{document}
\begin{tikzpicture}
  % [Format]{unterer Rand mm}{oberer Rand mm}
  % unterer Rand: ohne Legende 10; mit Legende \vfdbRandEintraege{<Zahl der Legendeneintraege>}
  % (dieselbe Zahl wie bei "vfdb legende" unten)
  \vfdbRahmen[standard]{\vfdbRandEintraege{2}}{3}
  \begin{axis}[vfdb=standard,          % standard | breit (zweifeldrig: siehe Anleitung)
      xmin=0, xmax=10, xtick distance=2,
      ymin=0, ymax=100, ytick distance=20,
      xlabel={Zeit $t$ in \si{\minute}},
      ylabel={Temperatur $T$ in \si{\degreeCelsius}},
      vfdb legende=2]                  % Zahl der Legendeneintraege (waehlt die Spalten)
    % Reihen erhalten Farbe und Linienstil automatisch in fester Reihenfolge
    \addplot table[x=t, y=A, col sep=semicolon] {
      t;A;B
      0;20;20
      5;60;45
      10;90;70
    };
    \addlegendentry{Messung A}
    \addplot table[x=t, y=B, col sep=semicolon] {
      t;A;B
      0;20;20
      5;60;45
      10;90;70
    };
    \addlegendentry{Messung B}
    % Daten aus einer Datei (CSV mit Dezimalpunkt, Semikolon getrennt):
    % \addplot table[x=t, y=A, col sep=semicolon] {daten.csv};
  \end{axis}
\end{tikzpicture}
\end{document}
